\section{E3 - LSTM/GRU trên chuỗi ảnh}

\subsection{Mục tiêu và các cách biến ảnh thành chuỗi}
E3 tiếp tục dùng Fashion-MNIST. Thay vì chỉ chọn một representation, báo cáo chuẩn bị ba cách để tạo đủ evidence cho comparison:
\begin{enumerate}
\item \textbf{Rows}: $T=28$, mỗi timestep có 28 feature.
\item \textbf{Columns}: $T=28$, mỗi timestep có 28 feature.
\item \textbf{Patch $4\times4$}: $T=49$, mỗi timestep có 16 feature.
\end{enumerate}

\subsection{Manual LSTM và GRU}
Repository có manual cell theo phương trình cổng để kiểm chứng kiến thức. \code{nn.LSTM}/\code{nn.GRU} vẫn được giữ làm reference vì đề cho phép, nhưng không thay thế phần manual.

\paragraph{LSTM.} Báo cáo cần trình bày input, forget, candidate và output gates; chỉ rõ cách cập nhật $c_t$ và $h_t$.
\paragraph{GRU.} Báo cáo cần trình bày update/reset gates và hidden candidate; giải thích vì sao GRU có ít state hơn LSTM.

\safeimage{e3/e3_recurrent_cells.png}{Sơ đồ manual LSTM/GRU cell và luồng ảnh $\rightarrow$ sequence $\rightarrow$ classifier.}

\subsection{Ma trận thí nghiệm}
\begin{table}[H]\centering
\caption{Ma trận thí nghiệm E3.}
\begin{tabularx}{0.92\textwidth}{@{}lllX@{}}
\toprule
ID & Recurrent model & Representation & Mục đích \\
\midrule
E3-L-R/E3-G-R & LSTM / GRU & Rows & Baseline sequence \\
E3-L-C/E3-G-C & LSTM / GRU & Columns & Sensitivity theo hướng quét \\
E3-L-P/E3-G-P & LSTM / GRU & Patch $4\times4$ & Local block sequence \\
E1 baseline & MLP / best CNN & 2D image & So sánh non-recurrent \\
\bottomrule
\end{tabularx}
\end{table}

\subsection{Kết quả}
\begin{table}[H]\centering
\caption{So sánh E3 với baseline E1.}
\begin{tabular}{@{}lllrrrr@{}}
\toprule
Model & Representation & Params & Val Acc & Test Acc & Time/epoch & Best epoch \\
\midrule
LSTM & Row & \metricpending & \metricpending & \metricpending & \metricpending & \metricpending \\
GRU & Row & \metricpending & \metricpending & \metricpending & \metricpending & \metricpending \\
LSTM & Column & \metricpending & \metricpending & \metricpending & \metricpending & \metricpending \\
GRU & Column & \metricpending & \metricpending & \metricpending & \metricpending & \metricpending \\
LSTM & Patch & \metricpending & \metricpending & \metricpending & \metricpending & \metricpending \\
GRU & Patch & \metricpending & \metricpending & \metricpending & \metricpending & \metricpending \\
Best E1 CNN & 2D image & \metricpending & \metricpending & \metricpending & \metricpending & \metricpending \\
\bottomrule
\end{tabular}
\end{table}
\safeimage{e3/e3_learning_comparison.png}{Learning curves và so sánh E3 với MLP/CNN baseline.}

\subsection{Thảo luận E3}
\TODO{PHÂN TÍCH LSTM VS GRU, ROW VS COLUMN VS PATCH, PARAMETER/TIME TRADE-OFF VÀ VÌ SAO MÔ HÌNH 2D CNN CÓ THỂ CÓ LỢI THẾ TRÊN ẢNH.}
